\documentclass{article}
\usepackage[equations]{propositions}

%% The nolabel style: an item with no label at all.
%%
%% counter=none suppresses the counter, and has done so since the beginning,
%% but not a name -- so on its own it leaves the label standing wherever a
%% name is inherited from the environment or supplied by a style.  nolabel is
%% counter=none, name={} and reset=false, which covers all three.

\begin{document}

\section{No label, and no number consumed}

\noindent\rule{\textwidth}{0.3pt}
\begin{prop}
  \item Numbered. \label{n:1}
  \item[style=nolabel] No label at all, and the count must not advance.
    \label{n:none}
  \item Numbered again. \label{n:2}
\end{prop}
\noindent\rule{\textwidth}{0.3pt}

The text of the unlabelled item must start where the other items' text starts:
suppressing the label empties the label box, it does not remove it.

References: [\ref{n:1}] [\ref{n:none}] [\ref{n:2}] --- expected (1), then
nothing at all, then (2).  An item with nothing to refer to takes no reference
format either: before that rule was mirrored from the display side, the middle
one printed \texttt{()}.

\section{A name that would otherwise be inherited}

\texttt{counter=none} suppresses the counter but not the name, so it is not
enough by itself in either list below.

\begin{prop}[name=Inherited]
  \item An item taking the inherited name.
  \item[counter=none] \texttt{counter=none}: still labelled \texttt{Inherited}.
  \item[style=nolabel] \texttt{nolabel}: no label.
\end{prop}

\begin{prop}[style=bullet]
  \item An ordinary bullet.
  \item[counter=none] \texttt{counter=none}: still bulleted, the bullet being
    a name and not a counter.
  \item[style=nolabel] \texttt{nolabel}: no bullet.
\end{prop}

\section{An unlabelled aside between two sub-lists}

This is what \texttt{reset = false} is in the style for.  The aside interrupts
a run of sub-items; the run must carry on across it, and must still be
referred to through the item it belongs to.

\begin{prop}
  \item An outer item. \label{n:outer}
  \begin{prop}
    \item Sub a. \label{n:a}
  \end{prop}
  \item[style=nolabel] Commentary on why the first sub-item is not enough,
    and we also need:
  \begin{prop}
    \item Sub b, continuing the run. \label{n:b}
  \end{prop}
\end{prop}

References: \ref{n:outer}, \ref{n:a}, \ref{n:b} --- the last must carry the
outer item's number and continue the letters, so it is that number followed by
\texttt{b}, not a bare \texttt{b}: an item with no reference of its own does
not become the parent of what follows.

\section{Sub-items, \texttt{inlineprop}, and a named list}

\begin{prop}
  \item Outer. \label{n:o2}
  \begin{prop}
    \item[style=nolabel] An unlabelled sub-item.
    \item A normal sub-item, which must be lettered \texttt{a}: the unlabelled
      one took no letter. \label{n:sa}
  \end{prop}
\end{prop}

References: \ref{n:o2}, \ref{n:sa}.

\begin{inlineprop}
  \item[style=nolabel] An unlabelled \texttt{inlineprop} item runs on with no
    label,
  \item and here is a second one.
\end{inlineprop}

\begin{prop}
  \item[Physicalism] A named item. \label{n:named}
  \item[style=nolabel] An unlabelled item in a list of named ones.
  \item[Dualism] Another named item. \label{n:named2}
\end{prop}

References: \ref{n:named}, \ref{n:named2}.

\section{A gloss survives}

\texttt{counter=none} with a gloss keeps the gloss, which is the one thing an
otherwise empty label still contributes.

\begin{prop}
  \item[counter=none, gloss=a gloss] An item whose label is only a gloss.
\end{prop}

\end{document}
